In this work, we introduced a novel strategy for representing real-time single-cell trajectories as information-rich, low-dimensional feature vectors. Our method takes advantage of probabilistic interpretations of optimal transport theory, along with the theory of reproducing kernel Hilbert spaces, to distill both trajectory and phenotypic information into a unified representation. This powerful approach enables us to use optimal transport trajectory modeling as a starting point for various downstream tasks in single-cell analysis, rather than as an endpoint in and of itself.

We paid particular attention to the applications of cell fate clustering and quantification of experimental perturbations, both of which have been subject of recent interest in the field \citep{langeCellRankDirectedSinglecell2022, peidliScPerturbHarmonizedSingleCell2023}. We demonstrated that clustering on the trajectory feature vectors produces cell subsets that are more consistent over time under the optimal transport maps. Furthermore, we showed that this may be particularly useful in biological systems with heterogeneous cell states, such as CD4 T-cell polarization. We also explored the connection between kernel mean embedding of trajectories and the maximum mean discrepancy, and proposed using the distance between trajectory vectors as a metric for the divergence between experimental responses. We demonstrated how using this approach was able to detect true-positive perturbations with higher efficiency than alternatives that do not take trajectories into account. Given the central role of time across biological processes, the trajectory featurization approach proposed in this paper has the potential to enable the study of single-cell responses at unprecedented temporal and cellular resolution.

Future work may seek to extend this approach from the discrete cell sets in optimal transport, to continuous gene expression space. This would enable integration with generative cell trajectory models such as TrajectoryNet \citep{tongTrajectoryNetDynamicOptimal2020} or PRESCIENT \citep{yeoGenerativeModelingSinglecell2021}. Another direction for further investigation is the extension of the featurization from discrete to continuous time. This could connect this work to the rich body of work on stochastic process and time-series modeling, including methods such as Gaussian processes.